\documentclass[10pt]{article}
\usepackage[margin=20mm]{geometry}
\usepackage{amsmath,amssymb,fontspec,unicode-math,xurl,hyperref}
\IfFontExistsTF{DejaVu Serif}{\setmainfont{DejaVu Serif}}{\setmainfont{Latin Modern Roman}}
\setmathfont{latinmodern-math.otf}
\setlength{\parindent}{0pt}\setlength{\parskip}{6pt}\sloppy
\title{Initial Scientific Review of
DRAG: Data Reconstruction Attack
using Guided Diffusion}
\author{AI-assisted technical review}
\date{}
\begin{document}
\maketitle
Abstract: The supplied evidence demonstrates reconstruction risks under an explicit white-box split-inference threat model. A Major defect occurs in the printed core sampling rule: its guidance sign and noise-scale convention disagree with the stated minimization objective and with the pinned implementation. The release uses a negative guidance direction and adds an already scaled displacement directly, providing substantial counterevidence against concluding that the reported experiments used the erroneous printed rule. A separate DINO-metric provenance concern remains Needs Evidence. No attack or author code was executed.\par
\section*{Identity, materials, and scope}
This review concerns paper 0127, I25-43496, by Wa-Kin Lei, Jun-Cheng Chen, and Shang-Tse Chen, in the supplied 19-page ICML 2025 PDF. Page locations are one-based PDF pages and match printed numbering. The original hash and inspected source hashes were checked against the acquired manifests. The pinned repository is `ntuaislab/DRAG`, commit `96788024346302e8481a56871d6c4b4b2c26ba63`.\par
The main method, threat model, quantitative tables, defense descriptions, and implementation appendices were read. Static source tracing covered the guided step, diffusion wrapper, optimizer-state extraction, metric feature selection, and selected defense code. Original inputs were preserved. The review performs only analytic and scalar calculations, without running reconstruction, author scripts, downloading models, or accessing private data.\par
\section*{Supported scope and experimental checks}
The attack assumes complete knowledge of the client encoder. The inverse-network extension additionally uses public training data; the paper separates this from private fine-tuning data and distinguishes DRAG from DRAG++. These assumptions matter: the evidence is not a demonstration of a query-free attack against an unknown private encoder.\par
Table 1 supports a substantial deepest-layer advantage under its reported metrics. For CLIP-ViT-B/16 at layer 12, DRAG reports MS-SSIM \(0.6735\) versus the best displayed alternative \(0.3884\), LPIPS \(0.1857\) versus \(0.5790\), and DINO similarity \(0.9331\) versus \(0.6524\). Shallow-layer superiority is not universal, and the paper acknowledges that. The \(T=250\), \(k=16\) schedule gives 4,000 refinement steps, consistent with Appendix A.3’s iteration count. That count does not imply faster wall-clock performance than every baseline; the execution table supplies separate times.\par
The evaluation uses ten images from each of three datasets, with sample indices listed in the release README. This is enough to establish examples of vulnerability but is a limited basis for population-level privacy rates. No confidence intervals, repeated seeds, or raw per-image results were independently reconstructed. No conclusion of differential-privacy failure is inferred from a defense that never claimed such a guarantee.\par
\section*{Problems}
Problem 1 (The printed guided step has the wrong minimization sign and inconsistent noise units).  Assessment: Major for the core algorithm as printed. Equation (1), p. 3, minimizes representation distance plus regularization. Equation (8), p. 4, defines \(g_t=\nabla_{x_t}L\) but adds its positive unit direction in the guided noise. Algorithm 1, p. 5, repeats that sign after positive Adam-style moment scaling. Moreover, equation (8) gives the replacement noise norm \(\sqrt n\,\sigma_t\), whereas equation (3) multiplies the noise by \(\sigma_t\) again.\par
\subsection*{Assumptions and decisive sign check}
Hold the selected guidance vector nonzero, use positive \(\sigma_t\), and consider a fixed denoising mean. At full guidance \(w=1\), the printed equation is \[\epsilon_t=r\,\frac{g_t}{\|g_t\|},\qquad r=\sqrt n\,\sigma_t.\] Its component along the loss gradient is positive. In a local coordinate where the guidance gradient is the objective gradient at the current mean, the linearized loss change from a positive displacement is positive. For the elementary smooth example \(L(u)=u^2/2\) at \(u=1\), a displacement \(+0.01\) increases the loss from \(0.5\) to \(0.51005\). A displacement \(-0.1\) decreases it to \(0.405\).\par
This scalar calculation does not claim that a nonconvex diffusion trajectory must decrease its loss at every step. It isolates the sign of the intended guidance component, which is the issue. It also does not rely on differentiating a discontinuous operation or on an actual image. On the first optimizer-state update in Algorithm 1, Adam’s positive normalization preserves each nonzero gradient sign; it does not supply the missing minus sign.\par
\subsection*{Noise-scale check}
Equation (3) writes the stochastic DDIM displacement as \(\sigma_t\epsilon_t\) for standard-normal \(\epsilon_t\). Equation (8), read literally as a replacement for that symbol, assigns it norm \(r=\sqrt n\,\sigma_t\). The resulting displacement consequently has norm \(\sqrt n\,\sigma_t^2\), rather than \(r\). For \(n=4\) and \(\sigma_t=0.1\), the two radii are \(0.02\) and \(0.2\). This affects every nonunit noise scale, independently of the sign issue.\par
The two consistent conventions are either a unit-variance noise variable subsequently multiplied by \(\sigma_t\), or an already scaled displacement added directly. The paper mixes those conventions. Simply changing the gradient sign without clarifying the stochastic variable leaves this inconsistency.\par
\subsection*{Pinned-source counterevidence}
The implementation does not use the printed sign and double scaling. In `run\_drag\_ldm.py`, lines 301–319, the scalar loss is differentiated, a normalized optimizer direction \(v\) is extracted, and the guided endpoint is explicitly `d\_star = -r * v`. The final displacement \(d\) is normalized to radius \(r\) and passed as `noise=d`. In `models/diffusion.py`, lines 81–98, supplied noise is added directly, without another \(\sigma_t\) multiplier. The pixel-space path in `run\_drag\_dm.py`, lines 280–298, uses the same convention.\par
This is strong evidence for a publication/implementation mismatch, not proof of failed historical reconstructions. The original experiments may have used the correct release behavior. The review does not replace that counterevidence with a claim that the reported images or metrics must be wrong.\par
\subsection*{Core dependency and correction bounds}
The affected rule is the central operation that makes diffusion conditional on the target representation. A reader implementing the stated equations obtains a different guidance direction and stochastic scale. This is a material core-algorithm error, assessed Major even though a notation/sign repair is short and other empirical results may survive.\par
A precise repair is to define the guided quantity as a displacement, use a negative loss-gradient direction, and add that displacement directly to the DDIM mean, consistently with the inspected release. Alternatively, retain the standardized-noise convention throughout. No author code was changed, no reconstruction was run, and no corrected score or convergence guarantee is claimed.\par
Problem 2 (The DINO feature used for evaluation needs provenance clarification).  Assessment: Needs Evidence. Section 4.3, p. 5, describes DINO ViT-S/16 image similarity as a semantic metric. The pinned `models/distance.py`, lines 20–57, returns `pooler\_output`, replaces the final layer normalization with the identity, and includes an explicit warning that pooler dense weights and bias were newly initialized rather than loaded from the checkpoint.\par
\subsection*{What is observed and what is not}
The provider’s model card describes the last hidden state of the CLS token as an image representation and gives an example reading `last\_hidden\_state`. It does not establish that the release’s extra pooled output is a pretrained semantic head. The release pins Transformers 4.50.3, but direct retrieval of that version’s implementation failed during this review. No checkpoint tensors were loaded, and the actual metric outputs used in the historical runs were unavailable.\par
The concern is scientifically meaningful because an unverified projection need not preserve cosine geometry. For example, \(u=(1,0)\) and \(v=(1,1)\) have cosine \(1/\sqrt2\), whereas applying \(\operatorname{diag}(10,1)\) to both gives \(10/\sqrt{101}\approx0.995\). This is an illustrative non-isometry, not an assertion that the released pooler equals that matrix or that its numerical errors have that size.\par
\subsection*{Counterevidence, impact, and needed resolution}
A fixed, shared projection can still preserve useful information, and using the same transform across methods can retain comparative value. In addition, MS-SSIM and LPIPS independently support many of the reported improvements. Neither the benchmark ordering nor the magnitude of DINO-score distortion has been demonstrated to change.\par
Resolution requires the actual feature-extraction definition, loaded-state key report, checkpoint identity, random initialization/seed information if applicable, and a comparison against the intended pretrained descriptor on the existing reconstructions. Until then, the DINO column’s interpretation and reproducibility remain uncertain. No claim that every metric or the privacy-risk conclusion fails is warranted.\par
\section*{Additional limits and overall assessment}
The token-permutation discussion is compatible with permutation-equivariant transformer blocks when shuffling occurs after positional embeddings and the CLS token is fixed. Dropping tokens is a separate operation that can affect utility. The adaptive token-reordering source maximizes sums of log probabilities through assignment, which matches the stated joint-probability interpretation; no issue was found in that scoped correspondence.\par
NoPeek is explicitly adapted from pretraining to fine-tuning and to the selected feature-distance setting in Appendix D. The reported defense comparison therefore concerns that implemented adaptation. The review does not establish failure of every possible NoPeek training procedure or future defense. The \(\ell_2\) penalty’s informal range language also should not be mistaken for a hard box constraint; the release separately clamps images before encoder evaluation.\par
The supported recommendation is Major for the printed guided sampler, with a distinct Needs Evidence metric concern. The historical empirical results remain unreplicated and are not disproved by these checks.\par
Scientific assessment | Major; adjudication pending\par
Repair verification | Sign/units reconciled analytically with release; no rerun\par
Author adoption | Not sought; unknown\par
Remaining work | Resolve metric provenance and the stated scientific evidence gaps\par
Independent calculations | Completed; exit code 0\par
Author-code execution | None\par
Focused acceptance | Recorded separately in provenance; adjudication pending\par
9 W.-K. Lei, J.-C. Chen, and S.-T. Chen, DRAG: Data Reconstruction Attack using Guided Diffusion, ICML 2025, https://proceedings.mlr.press/v267/lei25a.html; supplied final PDF. `ntuaislab/DRAG`, commit `96788024346302e8481a56871d6c4b4b2c26ba63`, files and lines cited above. Hugging Face model card for `facebook/dino-vits16`, consulted 24 October 2026, https://huggingface.co/facebook/dino-vits16.\par
\section*{Supporting checks}
The evidence includes source excerpts, independent scalar checks, the external metric-source record and visual inspection of source page 4.\par
The initial review chat performed all scientific analysis and writing. Auxiliary assistance was administrative only. Human reviews are empty, scientific decisions pending. Originals/manifests remain unchanged.\par
Scientific adjudication remains pending and human reviews remain empty.\par
\end{document}
